
\documentclass[hyperref={pdfpagelabels=false},ngerman]{beamer}

% stop font warning
\let\Tiny=\tiny
\providecommand\thispdfpagelabel[1]{}

\usepackage[english]{babel}
\usepackage{lmodern}
\usepackage[T1]{fontenc}
\usepackage[utf8]{inputenc}
\usepackage{graphicx,import}
\DeclareGraphicsRule{*}{mps}{*}{} 
\DeclareGraphicsExtensions{.pdf}
\usepackage{amsmath,amssymb,amstext,amsfonts} % mathrsfs
\usepackage{array,booktabs,tabularx}
\usepackage{tikz}
\usetikzlibrary{shapes,calc,arrows,positioning}
\tikzstyle{block} = [rectangle, draw, text width=7em, text centered, minimum height=2em]
\tikzstyle{arrow} = [draw, -latex, thick]
\tikzstyle{arrow2} = [draw, latex-latex, thick]
\tikzstyle{quark}  = [rectangle, draw, fill=yellow, minimum width=2em, text centered, minimum height=2em]
\tikzstyle{lepton} = [rectangle, draw, fill=red!50, minimum width=2em, text centered, minimum height=2em]
\tikzstyle{gauge}  = [circle   , draw, fill=green , minimum size=2em, inner sep=0pt, text centered]
\tikzstyle{scalar} = [diamond  , draw, fill=blue!40, minimum width=2.3em, text centered, minimum height=2.3em, inner sep=0pt]
\tikzstyle{goldstone} = [diamond, draw, dashed, fill=blue!30, minimum width=2.3em, text centered, minimum height=2.3em, inner sep=0pt]
\tikzstyle{squark}   = [diamond, draw, fill=yellow, minimum width=2.3em, text centered, minimum height=2.3em, inner sep=0pt]
\tikzstyle{slepton}  = [diamond, draw, fill=red!50, minimum width=2.3em, text centered, minimum height=2.3em, inner sep=0pt]
\tikzstyle{gaugino}  = [rectangle, draw, fill=green , minimum size=2em, inner sep=0pt, text centered]
\tikzstyle{higgsino} = [rectangle, draw, fill=blue!40  , minimum width=2em, text centered, minimum height=2em]
\tikzstyle{inert}    = [diamond  , draw, fill=teal!80, minimum width=2.3em, text centered, minimum height=2.3em, inner sep=0pt]
\tikzstyle{inertino} = [rectangle, draw, fill=teal!80, minimum width=2em, text centered, minimum height=2em]
\tikzstyle{phantom}  = [rectangle, minimum width=2em, text centered, minimum height=2em]
\usepackage{slashed}
\usepackage{fixltx2e} % textsubscript
\usepackage{multirow}
\usepackage{tcolorbox}
\usepackage{pifont}
\usepackage{xspace}
\usepackage{hyperref}
\hypersetup{colorlinks,linkcolor=,urlcolor=blue}
\usepackage{listings}
\lstset{breaklines=true,
  breakatwhitespace=true,
%  numbers=left,
  numberstyle=\tiny,
  stepnumber=1,
  basicstyle=\ttfamily\footnotesize,
  commentstyle=\ttfamily\color{gray},
  postbreak={\mbox{{$\hookrightarrow$}}\space\space},
  breakindent=10pt,
  breakautoindent=false,
  showspaces=false,
  showstringspaces=false,
  frame=single}

\definecolor{darkgreen}{RGB}{0,176,0}

\newcommand{\cmark}{\ding{51}}%
\newcommand{\xmark}{\ding{55}}%
\newcommand{\eh}[1]{\,\mathsf{#1}}
\newcommand{\ok}{\textcolor{darkgreen}{\cmark}}
\newcommand{\notok}{\textcolor{red}{\xmark}}
\newcommand{\maybe}{\textcolor{gray}{\cmark}}
\newcommand{\meh}{\textcolor{gray}{\textbf{\huge\lower.1em\hbox{-}}}}
\renewcommand{\emph}[1]{\textbf{\textcolor{darkblue}{#1}}}

% set look of slides
\usetheme{Madrid}
\useoutertheme{default}
\useinnertheme{circles}
\usecolortheme{default}
\beamertemplatenavigationsymbolsempty % keine Navigationselemente
\setbeamersize{text margin left = 1cm, text margin right = 1cm}

% define footer
\makeatletter
\setbeamertemplate{footline}
{
  \hfill\hbox{\insertframenumber{} / \inserttotalframenumber\hspace*{4pt}}%
  \vskip3pt%
}
\makeatother
\usecolortheme{tud}

\title{Es geht ums Prinzip: Was steckt hinter den Gleichungen der Physik?}

\author[Robert Harlander, Alexander Voigt]{Robert Harlander and Alexander Voigt}

\date{Sommerakademie der\\ Studienstiftung des deutschen Volkes\\[0.9em] Olang, 03.--14.09.2018}

\institute[Aachen]{RWTH Aachen}
\subject{Prinzip,Physik,Studienstiftung,deutschen,Volkes}
\keywords{Prinzip,Physik,Studienstiftung,deutschen,Volkes}

%%%%%%%%%%%%%%%%%%%%%%%%%%%%%%%%%%%%%%%%%%%%%%%%%%%%%%%%%%%%%%%%%%%%%%%%%%%%%

\begin{document}

%%%%%%%%%%%%%%%%%%%%%%%%%%%%%%%%%%%%%%%%
\begin{frame}[plain]
  \tikz [remember picture,overlay]
  \node at
    ([yshift=1.3cm,xshift=4cm]current page.south)
    {\includegraphics[height=2cm]{images/RWTH_Logo}};
  \titlepage  
\end{frame}

%%%%%%%%%%%%%%%%%%%%%%%%%%%%%%%%%%%%%%%%

\begin{frame}{Ziel des Workshops}
  \begin{block}{Was ist ein Prinzip?}
    Prinzip = Grundsatz oder Leitidee
  \end{block}
  \vspace*{2em}
  \emph{Ziel des Workshops:}
  \begin{itemize}
  \item Kennenlernen verschiedener Gestaltungs-, Erkenntnis- und
    Wirkprinzipien
  \item Ableiten von Konsequenzen bzgl.\ Natur, Forschung,
    Gesellschaft
  \item Bewerten der Prinzipien
  \end{itemize}
\end{frame}

%%%%%%%%%%%%%%%%%%%%%%%%%%%%%%%%%%%%%%%%

\begin{frame}{Themen}
  \emph{Übergeordnet}:\\
  Ontologie/Metaphysik\\[1em]
  \emph{Erkenntnisprinzipien:}\\
  Rahmen für Formulierung und Suche nach Erkenntnissen.\\[0.5em]
  \textit{Beispiele:}\\
  mathematische Beschreibung, innere Konsistenz, Einfachheit/Eleganz,
  Reduktionismus, Natürlichkeit, Korrespondenzprinzip,
  Renormierbarkeit,
  Anthropisches Prinzip, Symmetrieprinzip\\[1em]
  \emph{Wirkprinzipien:}\\
  Gesetzmäßigkeiten/Regeln, Kausalbeziehungen, Strukturen.\\[0.5em]
  \textit{Beispiele:}\\
  Prinzip der kleinsten Wirkung, äußere/innere Symmetrien, Kausalität,
  Lokalität, Unitarität, Relativitaetsprinzip, Äquivalenzprinzip
\end{frame}

%%%%%%%%%%%%%%%%%%%%%%%%%%%%%%%%%%%%%%%%

\newcommand{\thema}[1]{\textcolor{blue}{#1}}
\newcommand{\person}[1]{(#1)}

\begin{frame}{Zeitplan -- Woche 1 (03.--07.09.2018)}
  \begin{tabularx}{\textwidth}{ll}
    % \toprule
    Tag & Thema \\
    \midrule
    Mo & Vorstellungsrunde\\
       & \thema{Was sind Prinzipien?} (Alexey Schwarzmann)\\
    Di & \thema{Physik und Mathematik: innere Konsistenz, Logik}\\
       & (Meike Neuwohner)\\
       & \thema{Prinzipien vs. Axiome} (Johann Kipping und Doran Nettig)\\
       & \thema{Prinzip der kleinsten Wirkung} (Friederike Ihssen)\\
    Mi & \thema{Einfachheit und Eleganz} (Kim Wegner und Julian Schulz)\\
       & \thema{Ontologie} (Mirjam Trieb und Stefan Ringer)\\
       & \thema{Reduktionismus} (Kaja Jurak)\\
    Do & \thema{Symmetrie-Prinzip} (Jonas Peteranderl)\\
       & \thema{Äußere Symmetrien} (Jonas Lübken und Aaron Wiebel)\\
       & \thema{Relativitaetsprinzip} (Linus Kameni Monkam)\\
    Fr & \thema{Korrespondenzprinzip} (Ronan Schwarz)\\
       & \thema{Natürlichkeit} (Niklas Merk)\\
       & \thema{Unitarität} (Leon-Xinchen Yu)\\
    % \bottomrule
  \end{tabularx}
\end{frame}

\begin{frame}{Zeitplan -- Woche 2 (10.--14.09.2018)}
  \begin{tabularx}{\textwidth}{lX}
    % \toprule
    Tag & Thema \\
    \midrule
    Mo & \thema{Prinzipien in der Chemie} (Andreas Körner)\\
       & \thema{Anthropisches Prinzip} (Stefan Futterknecht)\\
       & \thema{Heisenberg und Philosophie} (Anna Swetlana Höfler)\\
    Di & Gruppenarbeit zu \textit{``Gibt es Mathematik auch ohne den Menschen?''} und \textit{``Lässt sich die Natur prinzipiell durch kleinste Entitäten beschreiben?''} \\
    Mi & Ausarbeitung Pro/Kontra\\
    Do & Podiumsdiskussionen (mit je einem Moderator)\\
    Fr & \thema{Kausalität und Lokalität} (Anna Geisler)\\
    % \bottomrule
  \end{tabularx}
\end{frame}

%%%%%%%%%%%%%%%%%%%%%%%%%%%%%%%%%%%%%%%%

\begin{frame}{Vorstellungsrunde}
  \begin{center}
    \includegraphics[width=0.3\textwidth]{images/haende_mann-frau}
  \end{center}
\end{frame}

\end{document}
